\QHeading{第1题　口算。}
\par\noindent


\begin{oralgrid}
\OralItem{1}{120\div6}{20}
\OralItem{2}{13+400}{413}
\OralItem{3}{30\times4}{120}
\OralItem{4}{58+12}{70}
\OralItem{5}{25\times2}{50}
\OralItem{6}{270\div90}{3}
\OralItem{7}{15\times30}{450}
\OralItem{8}{350\div7}{50}
\OralItem{9}{125\times8}{1000}
\OralItem{10}{360\div4}{90}
\OralItem{11}{60\times15}{900}
\OralItem{12}{450\div9}{50}\end{oralgrid}
\par

\QHeading{口算思路补充（对应上表题号）}

\Answer{上面列出得数，下面按小题给出不同计算思路。计数单位法适用于小数，整数题可比较凑整、拆分或商不变等方法。}
\MethodTitle{第1小题：$120\div 6$}

\MethodBlock{口算思路一：利用商不变的性质}{被除数和除数同时乘同一个非零数，商不变。}{\begin{aligned}[t]&\phantom{={}}120\div 6\\
&=(120\times 5)\div (6\times 5)\\
&=600\div 30\\
&=20\end{aligned}}

\MethodBlock{口算思路二：拆分除数，连续相除}{}{\begin{aligned}[t]&\phantom{={}}120\div 6\\
&=120\div (2\times 3)\\
&=120\div 2\div 3\\
&=60\div 3\\
&=20\end{aligned}}
\MethodTitle{第5小题：$25\times 2$}

\MethodBlock{口算思路一：按数位拆分计算}{把一个因数拆成整十、整百数与余下部分之和。}{\begin{aligned}[t]&\phantom{={}}25\times 2\\
&=2\times (20+5)\\
&=2\times 20+2\times 5\\
&=40+10\\
&=50\end{aligned}}
\MethodTitle{第6小题：$270\div 90$}

\MethodBlock{口算思路一：同时去掉相同个数的0}{}{\begin{aligned}[t]&\phantom{={}}270\div 90\\
&=(270\div 10)\div (90\div 10)\\
&=27\div 9\\
&=3\end{aligned}}

\MethodBlock{口算思路二：拆分除数，连续相除}{}{\begin{aligned}[t]&\phantom{={}}270\div 90\\
&=270\div (2\times 45)\\
&=270\div 2\div 45\\
&=135\div 45\\
&=3\end{aligned}}
\MethodTitle{第7小题：$15\times 30$}

\MethodBlock{口算思路一：按数位拆分计算}{把一个因数拆成整十、整百数与余下部分之和。}{\begin{aligned}[t]&\phantom{={}}15\times 30\\
&=30\times (10+5)\\
&=30\times 10+30\times 5\\
&=300+150\\
&=450\end{aligned}}
\MethodTitle{第9小题：$125\times 8$}

\MethodBlock{口算思路一：凑整后补偿}{把接近整十、整百的因数凑整，再补回或减去多算的部分。}{\begin{aligned}[t]&\phantom{={}}125\times 8\\
&=125\times (10-2)\\
&=125\times 10-125\times 2\\
&=1250-250\\
&=1000\end{aligned}}

\MethodBlock{口算思路二：按数位拆分计算}{把一个因数拆成整十、整百数与余下部分之和。}{\begin{aligned}[t]&\phantom{={}}125\times 8\\
&=8\times (100+25)\\
&=8\times 100+8\times 25\\
&=800+200\\
&=1000\end{aligned}}
\MethodTitle{第10小题：$360\div 4$}

\MethodBlock{口算思路一：利用商不变的性质}{被除数和除数同时乘同一个非零数，商不变。}{\begin{aligned}[t]&\phantom{={}}360\div 4\\
&=(360\times 5)\div (4\times 5)\\
&=1800\div 20\\
&=90\end{aligned}}

\MethodBlock{口算思路二：拆分除数，连续相除}{}{\begin{aligned}[t]&\phantom{={}}360\div 4\\
&=360\div (2\times 2)\\
&=360\div 2\div 2\\
&=180\div 2\\
&=90\end{aligned}}
\MethodTitle{第11小题：$60\times 15$}

\MethodBlock{口算思路一：按数位拆分计算}{把一个因数拆成整十、整百数与余下部分之和。}{\begin{aligned}[t]&\phantom{={}}60\times 15\\
&=60\times (10+5)\\
&=60\times 10+60\times 5\\
&=600+300\\
&=900\end{aligned}}
\MethodTitle{第12小题：$450\div 9$}

\MethodBlock{口算思路一：拆分除数，连续相除}{}{\begin{aligned}[t]&\phantom{={}}450\div 9\\
&=450\div (3\times 3)\\
&=450\div 3\div 3\\
&=150\div 3\\
&=50\end{aligned}}


\par\noindent

\QHeading{第2题　比一比，算一算。}

\CalcBlock{1}{\begin{aligned}[t]&\phantom{={}}\color{black}480\div12+8\times5\\&\color{answerred}=40+40\\&\color{answerred}=80\end{aligned}}
\par

\CalcBlock{2}{\begin{aligned}[t]&\phantom{={}}\color{black}480\div(12+8)\times5\\&\color{answerred}=480\div20\times5\\&\color{answerred}=24\times5\\&\color{answerred}=120\end{aligned}}
\CalcBlock[方法一]{3}{\begin{aligned}[t]&\phantom{={}}\color{black}480\div[(12+8)\times2]\\&\color{answerred}=480\div[20\times2]\\&\color{answerred}=480\div40\\&\color{answerred}=12\end{aligned}}
\MethodBlock{方法二：除以积，改成连续除}{连续除以两个数，等于除以这两个数的积。}{\begin{aligned}[t]&\phantom{={}}480\div [(12+8)\times 2]\\
&=480\div (12+8)\div 2\\
&=480\div 20\div 2\\
&=24\div 2\\
&=12\end{aligned}}

\CalcBlock[方法一]{4}{\begin{aligned}[t]&\phantom{={}}\color{black}200-50\times2+30\\&\color{answerred}=200-100+30\\&\color{answerred}=100+30\\&\color{answerred}=130\end{aligned}}
\MethodBlock{方法二：先处理括号或乘除，再交换、结合后先凑整}{保留每个数原来的加减号，先把能凑整的部分放在一起。}{\begin{aligned}[t]&\phantom{={}}200-50\times 2+30\\
&=200-100+30\\
&=(200+30)-100\\
&=230-100\\
&=130\end{aligned}}

\CalcBlock[方法一]{5}{\begin{aligned}[t]&\phantom{={}}\color{black}200-(50\times2+30)\\&\color{answerred}=200-(100+30)\\&\color{answerred}=200-130\\&\color{answerred}=70\end{aligned}}
\MethodBlock{方法二：先处理括号或乘除，再交换、结合后先凑整}{保留每个数原来的加减号，先把能凑整的部分放在一起。}{\begin{aligned}[t]&\phantom{={}}200-(50\times 2+30)\\
&=200-(100+30)\\
&=(200-30)-100\\
&=170-100\\
&=70\end{aligned}}

\MethodBlock{方法三：先处理括号或乘除，再去括号后计算}{加号后的括号去掉不变号，减号后的括号去掉要改变各项的加减号。}{\begin{aligned}[t]&\phantom{={}}200-(50\times 2+30)\\
&=200-(100+30)\\
&=200-100-30\\
&=100-30\\
&=70\end{aligned}}

\CalcBlock[方法一]{6}{\begin{aligned}[t]&\phantom{={}}\color{black}(200-50)\times(2+30)\\&\color{answerred}=150\times32\\&\color{answerred}=4800\end{aligned}}
\MethodBlock{方法二：乘法分配律}{把括号中的两部分分别相乘，再相加或相减。}{\begin{aligned}[t]&\phantom{={}}(200-50)\times (2+30)\\
&=(200-50)\times 2+(200-50)\times 30\\
&=150\times 2+150\times 30\\
&=300+4500\\
&=4800\end{aligned}}

\par\noindent

\QHeading{第3题　解决问题。}

\Problem{①体育用品商店每个足球售价61元，王老师带了500元，买了8个足球后，王老师还剩多少元？}
\par


\MethodTitle{方法一}
\Answer{\Formula{\begin{aligned}[t]&\phantom{={}}500-61\times8\\&=500-488\\&=12\text{（元）}\end{aligned}}\par 答：王老师还剩12元。}

\MethodTitle{方法二：先按整十单价算，再减去少算的费用}
\Answer{先把每个足球按60元算，实际每个还要多付1元。}
\Answer{\Formula{\begin{aligned}[t]&\phantom{={}}500-60\times 8-1\times 8\\
&=500-480-8\\
&=20-8\\
&=12\text{（元）}\end{aligned}}}
\Answer{答：王老师还剩12元。}
\Problem{②李师傅准备15天加工420个零件，实际提前3天完成了任务，实际每天多加工多少个零件？}

\MethodTitle{方法一}
\Answer{\Formula{\begin{aligned}[t]&\phantom{={}}420\div(15-3)-420\div15\\&=420\div12-420\div15\\&=35-28\\&=7\text{（个）}\end{aligned}}\par 答：实际每天多加工7个零件。}

\MethodTitle{方法二：比较实际这段时间内需要多做的零件}
\Answer{先求计划每天的加工量，再算实际12天内按计划本来能做多少，差额平均分到12天。}
\Answer{\textbf{计划每天加工}\par \Formula{\begin{aligned}[t]&\phantom{={}}420\div 15\\
&=28\text{（个）}\end{aligned}}}
\Answer{\textbf{实际加工天数}\par \Formula{\begin{aligned}[t]&\phantom{={}}15-3\\
&=12\text{（天）}\end{aligned}}}
\Answer{\textbf{12天内需要比原计划多加工}\par \Formula{\begin{aligned}[t]&\phantom{={}}420-28\times 12\\
&=420-336\\
&=84\text{（个）}\end{aligned}}}
\Answer{\textbf{实际每天多加工}\par \Formula{\begin{aligned}[t]&\phantom{={}}84\div 12\\
&=7\text{（个）}\end{aligned}}}
\Answer{答：实际每天多加工7个零件。}
